\documentclass[11pt]{article}

% ============================================================
% Page layout
% ============================================================

\usepackage[margin=1in]{geometry}
\usepackage{setspace}
\onehalfspacing
\usepackage{microtype}

% ============================================================
% Encoding and fonts
% ============================================================

\usepackage[T1]{fontenc}
\usepackage[utf8]{inputenc}
\usepackage{lmodern}

% ============================================================
% Math packages
% ============================================================

\usepackage{amsmath}
\usepackage{amssymb}
\usepackage{amsthm}
\usepackage{aliascnt}
\usepackage{mathtools}
\usepackage{bm}
\usepackage{bbm}

% ============================================================
% Tables and formatting
% ============================================================

\usepackage{booktabs}
\usepackage{array}
\usepackage{enumitem}
\usepackage{xcolor}

% ============================================================
% Citations, hyperlinks, and references
% ============================================================

\usepackage[round,authoryear]{natbib}

\usepackage[
    colorlinks=true,
    linkcolor=blue,
    citecolor=blue,
    urlcolor=blue
]{hyperref}

\usepackage[nameinlink,noabbrev]{cleveref}

% ============================================================
% Theorem environments
% ============================================================

\theoremstyle{plain}
\newtheorem{theorem}{Theorem}[section]

\newaliascnt{lemma}{theorem}
\newtheorem{lemma}[lemma]{Lemma}
\aliascntresetthe{lemma}

\newaliascnt{proposition}{theorem}
\newtheorem{proposition}[proposition]{Proposition}
\aliascntresetthe{proposition}

\newaliascnt{corollary}{theorem}
\newtheorem{corollary}[corollary]{Corollary}
\aliascntresetthe{corollary}

\theoremstyle{definition}
\newaliascnt{definition}{theorem}
\newtheorem{definition}[definition]{Definition}
\aliascntresetthe{definition}

\newaliascnt{assumption}{theorem}
\newtheorem{assumption}[assumption]{Assumption}
\aliascntresetthe{assumption}

\newaliascnt{example}{theorem}
\newtheorem{example}[example]{Example}
\aliascntresetthe{example}

\theoremstyle{remark}
\newaliascnt{remark}{theorem}
\newtheorem{remark}[remark]{Remark}
\aliascntresetthe{remark}

% Cleveref names for theorem-like environments with shared numbering.
\crefname{theorem}{Theorem}{Theorems}
\Crefname{theorem}{Theorem}{Theorems}
\crefname{lemma}{Lemma}{Lemmas}
\Crefname{lemma}{Lemma}{Lemmas}
\crefname{proposition}{Proposition}{Propositions}
\Crefname{proposition}{Proposition}{Propositions}
\crefname{corollary}{Corollary}{Corollaries}
\Crefname{corollary}{Corollary}{Corollaries}
\crefname{definition}{Definition}{Definitions}
\Crefname{definition}{Definition}{Definitions}
\crefname{assumption}{Assumption}{Assumptions}
\Crefname{assumption}{Assumption}{Assumptions}
\crefname{example}{Example}{Examples}
\Crefname{example}{Example}{Examples}
\crefname{remark}{Remark}{Remarks}
\Crefname{remark}{Remark}{Remarks}

% ============================================================
% Math operators
% ============================================================

\DeclareMathOperator{\E}{\mathbb{E}}
\DeclareMathOperator{\Var}{Var}
\DeclareMathOperator{\Cov}{Cov}
\DeclareMathOperator{\RMSPE}{RMSPE}
\DeclareMathOperator{\PH}{PH}
\DeclareMathOperator{\VR}{VR}
\DeclareMathOperator{\supp}{supp}
\DeclareMathOperator{\Law}{Law}

\DeclareMathOperator*{\argmin}{arg\,min}
\DeclareMathOperator*{\argmax}{arg\,max}

% ============================================================
% Common shortcuts
% ============================================================

\newcommand{\R}{\mathbb{R}}
\newcommand{\N}{\mathbb{N}}
\newcommand{\1}{\mathbbm{1}}

\newcommand{\SCM}{\mathrm{SCM}}
\newcommand{\TASCM}{\mathrm{TASCM}}

\newcommand{\pre}{\mathrm{pre}}
\newcommand{\post}{\mathrm{post}}
\newcommand{\train}{\mathrm{train}}
\newcommand{\val}{\mathrm{val}}

\newcommand{\topo}{\mathrm{top}}
\newcommand{\adm}{\mathrm{adm}}

\newcommand{\calB}{\mathcal B}
\newcommand{\calD}{\mathcal D}
\newcommand{\calF}{\mathcal F}
\newcommand{\calA}{\mathcal A}
\newcommand{\calR}{\mathcal R}
\newcommand{\calT}{\mathcal T}

\newcommand{\dB}{d_{\mathcal B}}
\newcommand{\dP}{d_{\mathcal P}}
\newcommand{\dH}{d_H}

\newcommand{\Ly}{L_y}
\newcommand{\Ltop}{L_{\mathrm{top}}}

\newcommand{\calLy}{\mathcal L_y}
\newcommand{\calLtop}{\mathcal L_{\mathrm{top}}}
\newcommand{\calQ}{\mathcal Q}

\newcommand{\tLy}{\widetilde L_y}
\newcommand{\tLtop}{\widetilde L_{\mathrm{top}}}

\newcommand{\tcalLy}{\widetilde{\mathcal L}_y}
\newcommand{\tcalLtop}{\widetilde{\mathcal L}_{\mathrm{top}}}
\newcommand{\tcalQ}{\widetilde{\mathcal Q}}

\newcommand{\tQ}{\widetilde Q}

\newcommand{\What}{\widehat w}
\newcommand{\Yhat}{\widehat Y}
\newcommand{\lhat}{\widehat\lambda}

% Probability convergence shortcut
\newcommand{\pto}{\xrightarrow{p}}

% ============================================================
% Document information
% ============================================================

\title{Theoretical Appendix for Topology-Augmented Synthetic Control}
\author{}
\date{}

% ============================================================
% Begin document
% ============================================================

\begin{document}

\maketitle

\appendix

% ============================================================
\section{Theoretical Appendix for Topology-Augmented Synthetic Control}
% ============================================================

This appendix provides the formal theoretical details for the
Topology-Augmented Synthetic Control Method (TASCM). The goal is to define
TASCM as a population M-estimator, establish consistency under uniform
convergence, derive an oracle risk bound, justify validation-based selection of
the topology penalty, and characterize a pairwise topology-dominance region.

Throughout this appendix, all claims are conditional on the stated assumptions.
The results do not imply that TASCM universally dominates conventional SCM.
Rather, they show that TASCM is a well-defined counterfactual-construction
estimator whose risk can be controlled under topology-relevant structural
conditions. The causal interpretation remains based on the usual SCM-style
conditions, such as no anticipation, no interference, and valid donor pools
\citep{abadie2010synthetic,abadie2021using}.

Unless otherwise stated, the asymptotic sequence is taken with
\[
T_0\to\infty,
\]
while \(J\), \(m\), and \(\ell\) are fixed. Allowing the donor dimension \(J\),
the embedding dimension \(m\), or the delay \(\ell\) to grow would require
additional complexity conditions and is outside the scope of this appendix.

% ------------------------------------------------------------
\subsection{Notation and Basic Setup}
% ------------------------------------------------------------

Consider a panel with \(J+1\) units and \(T\) periods. Unit \(i=1\) is treated
after period \(T_0\), while donor units \(j=2,\ldots,J+1\) remain untreated.
Potential outcomes are denoted by \(Y_{it}(0)\) and \(Y_{it}(1)\). The treatment
indicator is
\[
D_{it}
=
\1\{i=1,\ t>T_0\}.
\]
The observed outcome is
\[
Y_{it}
=
D_{it}Y_{it}(1)
+
(1-D_{it})Y_{it}(0).
\]
The treatment effect for the treated unit is
\[
\tau_{1t}
=
Y_{1t}(1)-Y_{1t}(0),
\qquad t>T_0.
\]
The primary estimand is the untreated post-treatment path
\[
\{Y_{1t}(0):t>T_0\}.
\]
Pointwise and average treatment effects are derived functionals of this
counterfactual path. Given an estimator \(\widehat Y_{1t}(0)\), the pointwise
treatment-effect estimator is
\[
\widehat \tau_{1t}
=
Y_{1t}-\widehat Y_{1t}(0).
\]

Let
\[
w=(w_2,\ldots,w_{J+1})'
\in
\Delta_J,
\]
where
\[
\Delta_J
=
\left\{
(w_2,\ldots,w_{J+1})\in\mathbb R^J:
w_j\ge 0,\ 
\sum_{j=2}^{J+1}w_j=1
\right\}.
\]
For \(w\in\Delta_J\), define the synthetic untreated path
\[
Y_{w,t}(0)
=
\sum_{j=2}^{J+1}w_jY_{jt}(0).
\]
Since donor units are untreated, \(Y_{jt}=Y_{jt}(0)\) for donor units under the
standard no-interference and no-exposure assumptions. The post-treatment
counterfactual estimator is
\[
\widehat Y_{1t}^{\TASCM}(0)
=
\sum_{j=2}^{J+1}\widehat w_jY_{jt},
\qquad t>T_0.
\]

Unless explicitly labeled as validation loss or post-treatment risk, all
empirical losses below are computed on the pre-treatment period.

% ------------------------------------------------------------
\subsection{Delay-Window Measures and Population Topological Objects}
% ------------------------------------------------------------

The first issue is to define the population counterpart of the finite-sample
topology loss. A raw persistence diagram computed from a finite pre-treatment
trajectory is not, by itself, a population object. We therefore define topology
through the distribution of standardized delay windows, following the
statistical TDA perspective that topological summaries estimate features of an
underlying geometric or distributional object \citep{chazal2021introduction}.

For the topological comparison, each finite trajectory is standardized before
delay embedding. For any finite path \(Y=(Y_1,\ldots,Y_n)\), define
\[
\bar Y_n
=
\frac{1}{n}\sum_{t=1}^nY_t,
\qquad
s_n(Y)
=
\left\{
\frac{1}{n}\sum_{t=1}^n(Y_t-\bar Y_n)^2
\right\}^{1/2}.
\]
Let
\[
S_n(Y)_t
=
\frac{Y_t-\bar Y_n}{s_n(Y)\vee \epsilon_s},
\qquad t=1,\ldots,n,
\]
where \(\epsilon_s>0\) is a fixed small constant preventing division by zero.
The topology loss below is computed from \(S_n(Y)\), not from the raw level
path. Thus the outcome loss \(L_y\) controls level fit, while the topology loss
compares standardized dynamic shape.

Fix an embedding dimension \(m\ge2\) and a delay \(\ell\ge1\). For each unit
\(i\), let
\[
Y_i^{\pre}
=
(Y_{i1},\ldots,Y_{iT_0})
\]
denote its pre-treatment path. Define the standardized delay vector
\[
X_{i,t}^{S,m,\ell}
=
\bigl(
S_{T_0}(Y_i^{\pre})_t,
S_{T_0}(Y_i^{\pre})_{t-\ell},
\ldots,
S_{T_0}(Y_i^{\pre})_{t-(m-1)\ell}
\bigr)
\in\mathbb R^m,
\]
for
\[
t=(m-1)\ell+1,\ldots,T_0.
\]

For a synthetic control \(w\), first form the synthetic pre-treatment path
\[
Y_w^{\pre}
=
\sum_{j=2}^{J+1}w_jY_j^{\pre}.
\]
Then define the standardized synthetic delay vector
\[
X_{w,t}^{S,m,\ell}
=
\bigl(
S_{T_0}(Y_w^{\pre})_t,
S_{T_0}(Y_w^{\pre})_{t-\ell},
\ldots,
S_{T_0}(Y_w^{\pre})_{t-(m-1)\ell}
\bigr)
\in\mathbb R^m.
\]
Thus the synthetic trajectory is constructed first, and only then standardized
and embedded. In general,
\[
X_{w,t}^{S,m,\ell}
\neq
\sum_{j=2}^{J+1}w_jX_{j,t}^{S,m,\ell},
\]
because the standardization operator is nonlinear. This is exactly the
implementation logic: TASCM computes topology on the standardized synthetic
path \(S_{T_0}(Y_w^{\pre})\), not on a weighted average of donor topological
features.

Let
\[
N_0=T_0-(m-1)\ell.
\]
Define the empirical standardized delay-window measure
\[
\widehat P_{w,T_0}^{S}
=
\frac{1}{N_0}
\sum_{t=(m-1)\ell+1}^{T_0}
\delta_{X_{w,t}^{S,m,\ell}}.
\]
For the treated unit,
\[
\widehat P_{1,T_0}^{S}
=
\frac{1}{N_0}
\sum_{t=(m-1)\ell+1}^{T_0}
\delta_{X_{1,t}^{S,m,\ell}}.
\]

For each \(w\in\Delta_J\), let \(P_w^{S}\) denote the population law associated
with the standardized delay vectors of the synthetic path
\(S_{T_0}(Y_w^{\pre})\). Similarly, let \(P_1^{S}\) denote the corresponding
population law for the treated unit, and let \(P_j^{S}\) denote the
standardized delay-window law for donor \(j\). The empirical measures
\(\widehat P_{w,T_0}^{S}\) and \(\widehat P_{1,T_0}^{S}\) are sample analogues
of these population laws. Uniform convergence of these empirical laws is stated
formally in Assumption A6 below.

Persistent homology is nonlinear and does not commute with convex averaging.
Therefore, TASCM first constructs the synthetic trajectory and then computes its
topological summary. In general,
\[
\Psi_1(P_w^{S})
=
\Phi\left\{
D_1\left(
\mathcal F(P_w^{S})
\right)
\right\}
\neq
\sum_{j=2}^{J+1}w_j
\Phi\left\{
D_1\left(
\mathcal F(P_j^{S})
\right)
\right\}
=
\sum_{j=2}^{J+1}w_j\Psi_1(P_j^{S}).
\]
This distinction is essential because the topology of a weighted synthetic
trajectory can differ qualitatively from a weighted average of the topological
summaries of its donor components. Thus the topology term compares the
persistent-homology-based summary of the standardized synthetic trajectory, not
a weighted average of donor persistent-homology summaries.

The population topological object is defined by a stable filtration operator.
Let \(\mathcal F\) be a filtration-generating operator that maps a probability
law \(P\) on \(\mathbb R^m\) into a tame filtration \(\mathcal F(P)\). Examples
include
\[
\mathcal F(P)
=
\text{DTM filtration generated by }P,
\]
\[
\mathcal F(P)
=
\text{kernel-distance filtration generated by }P,
\]
or, under stronger support-convergence assumptions,
\[
\mathcal F(P)
=
\text{Vietoris--Rips filtration on } \supp(P).
\]
DTM filtrations are natural population-level examples because
distance-to-measure functions are designed for geometric inference from
probability measures \citep{chazal2011geometric}.

Let \(D_1(\mathcal F(P))\) denote the one-dimensional persistence diagram of
the filtration \(\mathcal F(P)\). Let \(\mathcal D_1\) denote the space of
one-dimensional persistence diagrams equipped with a diagram metric such as the
bottleneck or \(p\)-Wasserstein metric. Let
\[
\Phi:\mathcal D_1\to \mathcal B
\]
be a stable representation of persistence diagrams into a Banach, Hilbert, or
finite-dimensional Euclidean space \(\mathcal B\). Examples include
persistence landscapes, persistence images, silhouettes, and stable kernel
embeddings. Persistence landscapes provide Banach-space summaries suitable for
statistical analysis \citep{bubenik2015statistical}, while persistence images
provide stable finite-dimensional vector representations
\citep{adams2017persistence}.

Define
\[
\Psi_1(P)
=
\Phi\{D_1(\mathcal F(P))\}.
\]
The population topology loss is
\[
\mathcal L_{\mathrm{top}}(w)
=
d_{\mathcal B}
\left(
\Psi_1(P_1^{S}),
\Psi_1(P_w^{S})
\right).
\]
The sample topology loss is
\[
L_{\mathrm{top},T_0}(w)
=
d_{\mathcal B}
\left(
\Psi_1(\widehat P_{1,T_0}^{S}),
\Psi_1(\widehat P_{w,T_0}^{S})
\right).
\]

\paragraph{Comment.}
The operator \(\mathcal F\) is kept generic in the main theory. DTM or
kernel-distance filtrations provide the cleanest distribution-level population
objects. Vietoris--Rips persistence on finite standardized delay-embedded point
clouds can be treated as a finite-sample implementation under stronger
Hausdorff/support-convergence conditions. This distinction prevents a mismatch
between population theory and empirical implementation. For computational
foundations of persistent homology, see \citet{zomorodian2005computing}.

% ------------------------------------------------------------
\subsection{Outcome Losses and Standardized TASCM Objective}
% ------------------------------------------------------------

Define the population outcome loss as the long-run pre-treatment risk
\[
\mathcal L_y(w)
=
\lim_{T_0\to\infty}
\frac{1}{T_0}
\sum_{t=1}^{T_0}
\left(
Y_{1t}(0)
-
\sum_{j=2}^{J+1}w_jY_{jt}(0)
\right)^2,
\]
whenever the limit exists. Under stationarity and ergodicity, this limit can
also be written as
\[
\mathcal L_y(w)
=
\mathbb E
\left[
\left(
Y_{1t}(0)
-
\sum_{j=2}^{J+1}w_jY_{jt}(0)
\right)^2
\right].
\]
The corresponding empirical pre-treatment loss is
\[
L_{y,T_0}(w)
=
\frac{1}{T_0}
\sum_{t=1}^{T_0}
\left(
Y_{1t}
-
\sum_{j=2}^{J+1}w_jY_{jt}
\right)^2.
\]

Let the population SCM target be
\[
w_{\SCM}^0
\in
\argmin_{w\in\Delta_J}\mathcal L_y(w).
\]
Choose fixed constants \(a,b>0\) to avoid division by zero and define
\[
s_y
=
\mathcal L_y(w_{\SCM}^0)\vee a,
\qquad
s_{\mathrm{top}}
=
\mathcal L_{\mathrm{top}}(w_{\SCM}^0)\vee b.
\]
The constants \(a\) and \(b\) are fixed normalization floors chosen before
estimation and do not use post-treatment outcomes.

The standardized population losses are
\[
\widetilde{\mathcal L}_y(w)
=
\frac{\mathcal L_y(w)}{s_y},
\]
and
\[
\widetilde{\mathcal L}_{\mathrm{top}}(w)
=
\frac{\mathcal L_{\mathrm{top}}(w)}{s_{\mathrm{top}}}.
\]
The standardized population TASCM objective is
\[
\widetilde{\mathcal Q}_{\lambda}(w)
=
\widetilde{\mathcal L}_y(w)
+
\lambda
\widetilde{\mathcal L}_{\mathrm{top}}(w).
\]
The population TASCM target is
\[
w_{\lambda}^{\star}
\in
\argmin_{w\in\Delta_J}
\widetilde{\mathcal Q}_{\lambda}(w).
\]

For the empirical objective, let
\[
\widehat w_{\SCM}
\in
\argmin_{w\in\Delta_J}L_{y,T_0}(w),
\]
where a deterministic tie-breaking rule is imposed if the minimizer is not
unique. Define sample scales
\[
\widehat s_y
=
L_{y,T_0}(\widehat w_{\SCM})\vee a,
\qquad
\widehat s_{\mathrm{top}}
=
L_{\mathrm{top},T_0}(\widehat w_{\SCM})\vee b.
\]
Then
\[
\widetilde L_{y,T_0}(w)
=
\frac{L_{y,T_0}(w)}{\widehat s_y},
\]
and
\[
\widetilde L_{\mathrm{top},T_0}(w)
=
\frac{L_{\mathrm{top},T_0}(w)}{\widehat s_{\mathrm{top}}}.
\]
The empirical TASCM objective is
\[
\widetilde Q_{T_0,\lambda}(w)
=
\widetilde L_{y,T_0}(w)
+
\lambda
\widetilde L_{\mathrm{top},T_0}(w).
\]

Because the topology term is generally nonlinear and nonconvex in \(w\), the
theoretical estimator is allowed to be an approximate global minimizer. We say
that \(\widehat w_\lambda\) is a \(\zeta_T\)-approximate TASCM minimizer if
\[
\widetilde Q_{T_0,\lambda}(\widehat w_\lambda)
\le
\inf_{w\in\Delta_J}
\widetilde Q_{T_0,\lambda}(w)
+
\zeta_T,
\]
where
\[
\zeta_T\pto 0.
\]
The exact global-minimizer case corresponds to \(\zeta_T=0\).

The use of objective-standardized losses is important. Without this
normalization, \(\lambda\) absorbs arbitrary unit differences between outcome
mean-squared error and topological distance. This objective standardization is
distinct from the trajectory standardization \(S_n(\cdot)\) used before
computing persistent homology.

% ------------------------------------------------------------
\subsection{Assumptions}
% ------------------------------------------------------------

\paragraph{Assumption A1: No anticipation.}
For \(t\le T_0\),
\[
Y_{1t}=Y_{1t}(0).
\]

\paragraph{Assumption A2: No interference and valid donors.}
For all donor units \(j=2,\ldots,J+1\) and all \(t\),
\[
Y_{jt}=Y_{jt}(0).
\]
Donor units are not affected by the treatment assigned to unit 1 and are not
subject to the same intervention.

\paragraph{Assumption A3: Compactness.}
The donor weight space \(\Delta_J\) is compact.

\paragraph{Assumption A4: Boundedness.}
There exists \(M<\infty\) such that
\[
|Y_{it}(0)|\le M
\]
for all \(i,t\).

\paragraph{Assumption A5: Stable and finite topological representation.}
There exists a metric \(d_{\mathcal P}\) on standardized delay-window laws and
a constant \(C_\Psi<\infty\) such that
\[
d_{\mathcal B}
\left(
\Psi_1(P),
\Psi_1(Q)
\right)
\le
C_\Psi d_{\mathcal P}(P,Q)
\]
for all relevant standardized delay-window laws \(P,Q\). In addition,
\[
\sup_{w\in\Delta_J}\mathcal L_{\mathrm{top}}(w)<\infty.
\]

For example, if \(\mathcal F\) is a DTM filtration, stability can be obtained
from Wasserstein stability of the DTM. If \(\Phi\) is a persistence landscape or
persistence image, additional Lipschitz stability follows from known stability
results for those representations
\citep{cohensteiner2007stability,chazal2016structure,bubenik2015statistical,adams2017persistence}.
For VR-on-support, \(d_{\mathcal P}(P,Q)\) should be interpreted as a support
metric such as
\[
d_H(\supp(P),\supp(Q)),
\]
rather than a weak-distribution metric.

\paragraph{Assumption A6: Uniform convergence of standardized delay-window laws.}
The treated standardized delay-window law and the synthetic standardized
delay-window laws satisfy
\[
d_{\mathcal P}
\left(
\widehat P_{1,T_0}^{S},
P_1^{S}
\right)
\pto 0,
\]
and
\[
\sup_{w\in\Delta_J}
d_{\mathcal P}
\left(
\widehat P_{w,T_0}^{S},
P_w^{S}
\right)
\pto 0.
\]
Because delay windows overlap, the empirical measures
\(\widehat P_{w,T_0}^{S}\) are based on dependent observations. This assumption
should therefore be interpreted as a stationarity, ergodicity, or mixing-type
condition for the standardized delay-window process, uniformly over
\(w\in\Delta_J\).

\paragraph{Assumption A7: Uniform convergence of outcome loss.}
\[
\sup_{w\in\Delta_J}
\left|
L_{y,T_0}(w)-\mathcal L_y(w)
\right|
\pto 0.
\]

\paragraph{Assumption A8: Scale consistency.}
The empirical normalizing constants satisfy
\[
\widehat s_y \pto s_y,
\qquad
\widehat s_{\mathrm{top}}\pto s_{\mathrm{top}}.
\]
If the SCM minimizer is not unique, the deterministic tie-breaking rule above
makes \(\widehat w_{\SCM}\) and hence \(\widehat s_{\mathrm{top}}\) well-defined.
This condition is automatic if the population SCM minimizer is unique and the
topology loss is continuous at that minimizer.

\paragraph{Assumption A9: Continuity and unique population minimizer.}
For a fixed \(\lambda\ge0\), the population objective
\[
\widetilde{\mathcal Q}_{\lambda}(w)
\]
is continuous on \(\Delta_J\) and has a unique minimizer \(w_\lambda^\star\)
over \(\Delta_J\). Continuity follows, for example, if \(w\mapsto P_w^{S}\) is
continuous under \(d_{\mathcal P}\) and \(\Psi_1\) is Lipschitz as in
Assumption A5.

\paragraph{Assumption A10: Topology-augmented pre-post stability.}
Define the post-treatment counterfactual mean-squared risk
\[
\mathcal R(w)
=
\frac{1}{T_1}
\sum_{t=T_0+1}^{T}
\left[
Y_{1t}(0)
-
\sum_{j=2}^{J+1}w_jY_{jt}(0)
\right]^2,
\qquad
T_1=T-T_0.
\]
The risk \(\mathcal R(w)\) is infeasible because \(Y_{1t}(0)\) is unobserved for
\(t>T_0\). It is used only for theoretical evaluation.

There exist constants \(C_y,C_{\mathrm{top}}<\infty\) and a remainder
\(r_T\to0\) such that for all \(w\in\Delta_J\),
\[
\mathcal R(w)
\le
C_y\widetilde{\mathcal L}_y(w)
+
C_{\mathrm{top}}\widetilde{\mathcal L}_{\mathrm{top}}(w)
+
r_T.
\]
This is a structural bridge from pre-treatment balance to post-treatment
counterfactual validity. It is not implied by persistent homology alone. It
plays a role analogous to invariance or time-invariant factor structure in
conventional synthetic-control designs.

\paragraph{Assumption A11: Approximate topology-augmented synthetic representation.}
There exists a deterministic or population-defined sequence
\(w_T^\circ\in\Delta_J\) such that
\[
\widetilde{\mathcal L}_y(w_T^\circ)
\le
\epsilon_{y,T},
\qquad
\widetilde{\mathcal L}_{\mathrm{top}}(w_T^\circ)
\le
\epsilon_{\mathrm{top},T},
\]
where
\[
\epsilon_{y,T}\to0,
\qquad
\epsilon_{\mathrm{top},T}\to0.
\]
This assumption is stronger than the usual SCM convex-hull condition. It
requires the treated unit to be approximable not only in level outcomes but
also in the topological summary of the standardized synthetic trajectory. Since
persistent homology is nonlinear, level balance does not imply topology
balance.

% ------------------------------------------------------------
\subsection{A Sufficient Structure for Pre--Post Stability}
% ------------------------------------------------------------

The pre-post stability condition in Assumption A10 is high-level. The following
proposition gives one sufficient structure under which such a condition can be
motivated.

\begin{proposition}[A sufficient structure for Assumption A10]
\label{prop:prepost-structure}
Suppose untreated outcomes admit the decomposition
\[
Y_{it}(0)
=
\mu_i'f_t
+
\eta_i'g_t
+
\varepsilon_{it},
\]
where \(f_t\) are level factors, \(g_t\) are dynamic factors, and
\(\varepsilon_{it}\) is a mean-zero approximation error with uniformly bounded
second moment.

Suppose post-treatment factor moments are bounded:
\[
\frac{1}{T_1}
\sum_{t=T_0+1}^{T}
\|f_t\|^2
\le
K_f,
\qquad
\frac{1}{T_1}
\sum_{t=T_0+1}^{T}
\|g_t\|^2
\le
K_g.
\]
Suppose further that there exist constants \(K_y,K_{\mathrm{top}}<\infty\) such
that, for all \(w\in\Delta_J\),
\[
\left\|
\mu_1-\sum_{j=2}^{J+1}w_j\mu_j
\right\|^2
\le
K_y\widetilde{\mathcal L}_y(w),
\]
and
\[
\left\|
\eta_1-\sum_{j=2}^{J+1}w_j\eta_j
\right\|^2
\le
K_{\mathrm{top}}
\widetilde{\mathcal L}_{\mathrm{top}}(w).
\]

Finally, define the post-treatment approximation residual
\[
u_{w,t}
=
\varepsilon_{1t}
-
\sum_{j=2}^{J+1}w_j\varepsilon_{jt}.
\]
Assume that this residual satisfies the uniform vanishing condition
\[
\sup_{w\in\Delta_J}
\frac{1}{T_1}
\sum_{t=T_0+1}^{T}
u_{w,t}^2
=
o(1).
\]
Then there exist constants \(C_y,C_{\mathrm{top}}<\infty\) and a remainder
\(r_T\to0\) such that
\[
\mathcal R(w)
\le
C_y\widetilde{\mathcal L}_y(w)
+
C_{\mathrm{top}}\widetilde{\mathcal L}_{\mathrm{top}}(w)
+
r_T.
\]
\end{proposition}

\begin{proof}
For any \(w\in\Delta_J\),
\[
Y_{1t}(0)-Y_{w,t}(0)
=
\left(
\mu_1-\sum_{j=2}^{J+1}w_j\mu_j
\right)'f_t
+
\left(
\eta_1-\sum_{j=2}^{J+1}w_j\eta_j
\right)'g_t
+
u_{w,t}.
\]
Using
\[
(a+b+c)^2\le 3a^2+3b^2+3c^2,
\]
we obtain
\[
\begin{aligned}
\mathcal R(w)
&=
\frac{1}{T_1}
\sum_{t=T_0+1}^{T}
\left[
Y_{1t}(0)-Y_{w,t}(0)
\right]^2
\\
&\le
3
\left\|
\mu_1-\sum_{j=2}^{J+1}w_j\mu_j
\right\|^2
\frac{1}{T_1}
\sum_{t=T_0+1}^{T}
\|f_t\|^2
\\
&\quad+
3
\left\|
\eta_1-\sum_{j=2}^{J+1}w_j\eta_j
\right\|^2
\frac{1}{T_1}
\sum_{t=T_0+1}^{T}
\|g_t\|^2
\\
&\quad+
3
\frac{1}{T_1}
\sum_{t=T_0+1}^{T}
u_{w,t}^2.
\end{aligned}
\]
By the bounded post-treatment factor moment assumptions,
\[
\frac{1}{T_1}
\sum_{t=T_0+1}^{T}
\|f_t\|^2
\le K_f,
\qquad
\frac{1}{T_1}
\sum_{t=T_0+1}^{T}
\|g_t\|^2
\le K_g.
\]
Therefore,
\[
\begin{aligned}
\mathcal R(w)
&\le
3K_f
\left\|
\mu_1-\sum_{j=2}^{J+1}w_j\mu_j
\right\|^2
+
3K_g
\left\|
\eta_1-\sum_{j=2}^{J+1}w_j\eta_j
\right\|^2
\\
&\quad+
3
\frac{1}{T_1}
\sum_{t=T_0+1}^{T}
u_{w,t}^2.
\end{aligned}
\]
Using the assumed loading-control inequalities,
\[
\left\|
\mu_1-\sum_{j=2}^{J+1}w_j\mu_j
\right\|^2
\le
K_y\widetilde{\mathcal L}_y(w),
\]
and
\[
\left\|
\eta_1-\sum_{j=2}^{J+1}w_j\eta_j
\right\|^2
\le
K_{\mathrm{top}}
\widetilde{\mathcal L}_{\mathrm{top}}(w),
\]
we get
\[
\mathcal R(w)
\le
3K_fK_y
\widetilde{\mathcal L}_y(w)
+
3K_gK_{\mathrm{top}}
\widetilde{\mathcal L}_{\mathrm{top}}(w)
+
3
\frac{1}{T_1}
\sum_{t=T_0+1}^{T}
u_{w,t}^2.
\]
Define
\[
C_y=3K_fK_y,
\qquad
C_{\mathrm{top}}=3K_gK_{\mathrm{top}},
\]
and
\[
r_T
=
3
\sup_{w\in\Delta_J}
\frac{1}{T_1}
\sum_{t=T_0+1}^{T}
u_{w,t}^2.
\]
By the uniform vanishing residual condition,
\[
r_T\to0.
\]
Since
\[
3
\frac{1}{T_1}
\sum_{t=T_0+1}^{T}
u_{w,t}^2
\le
r_T
\]
for every \(w\in\Delta_J\), it follows that
\[
\mathcal R(w)
\le
C_y\widetilde{\mathcal L}_y(w)
+
C_{\mathrm{top}}\widetilde{\mathcal L}_{\mathrm{top}}(w)
+
r_T.
\]
This proves the claim.
\end{proof}

\begin{remark}
This proposition is not meant to claim that topology balance is always useful.
It states one sufficient structure: topology helps when the topological loss
controls a dynamic component of the untreated path that is relevant for
post-treatment counterfactual risk and not fully captured by level fit.
\end{remark}

% ------------------------------------------------------------
\subsection{Lemma 1: Uniform Convergence of the Topology Loss}
% ------------------------------------------------------------

\begin{lemma}[Uniform convergence of topology loss]
\label{lem:topology-uniform}
Under Assumptions A5 and A6,
\[
\sup_{w\in\Delta_J}
\left|
L_{\mathrm{top},T_0}(w)
-
\mathcal L_{\mathrm{top}}(w)
\right|
\pto 0.
\]
\end{lemma}

\begin{proof}
Recall that
\[
L_{\mathrm{top},T_0}(w)
=
d_{\mathcal B}
\left(
\Psi_1(\widehat P_{1,T_0}^{S}),
\Psi_1(\widehat P_{w,T_0}^{S})
\right)
\]
and
\[
\mathcal L_{\mathrm{top}}(w)
=
d_{\mathcal B}
\left(
\Psi_1(P_1^{S}),
\Psi_1(P_w^{S})
\right).
\]
By the reverse triangle inequality in the metric space
\((\mathcal B,d_{\mathcal B})\),
\[
\begin{aligned}
&
\left|
L_{\mathrm{top},T_0}(w)
-
\mathcal L_{\mathrm{top}}(w)
\right|
\\
&\le
d_{\mathcal B}
\left(
\Psi_1(\widehat P_{1,T_0}^{S}),
\Psi_1(P_1^{S})
\right)
+
d_{\mathcal B}
\left(
\Psi_1(\widehat P_{w,T_0}^{S}),
\Psi_1(P_w^{S})
\right).
\end{aligned}
\]
By Assumption A5,
\[
d_{\mathcal B}
\left(
\Psi_1(\widehat P_{1,T_0}^{S}),
\Psi_1(P_1^{S})
\right)
\le
C_\Psi d_{\mathcal P}(\widehat P_{1,T_0}^{S},P_1^{S}),
\]
and
\[
d_{\mathcal B}
\left(
\Psi_1(\widehat P_{w,T_0}^{S}),
\Psi_1(P_w^{S})
\right)
\le
C_\Psi d_{\mathcal P}(\widehat P_{w,T_0}^{S},P_w^{S}).
\]
Therefore,
\[
\sup_{w\in\Delta_J}
\left|
L_{\mathrm{top},T_0}(w)
-
\mathcal L_{\mathrm{top}}(w)
\right|
\le
C_\Psi d_{\mathcal P}(\widehat P_{1,T_0}^{S},P_1^{S})
+
C_\Psi
\sup_{w\in\Delta_J}
d_{\mathcal P}(\widehat P_{w,T_0}^{S},P_w^{S}).
\]
The right-hand side converges to zero in probability by Assumption A6. Hence
\[
\sup_{w\in\Delta_J}
\left|
L_{\mathrm{top},T_0}(w)
-
\mathcal L_{\mathrm{top}}(w)
\right|
\pto 0.
\]
\end{proof}

% ------------------------------------------------------------
\subsection{Lemma 2: Uniform Convergence of the Standardized Objective}
% ------------------------------------------------------------

\begin{lemma}[Uniform convergence of standardized TASCM objective]
\label{lem:objective-uniform}
Suppose Assumptions A4--A8 hold. Then, for fixed \(\lambda<\infty\),
\[
\sup_{w\in\Delta_J}
\left|
\widetilde Q_{T_0,\lambda}(w)
-
\widetilde{\mathcal Q}_{\lambda}(w)
\right|
\pto 0.
\]
\end{lemma}

\begin{proof}
By Assumption A7,
\[
\sup_{w\in\Delta_J}
|L_{y,T_0}(w)-\mathcal L_y(w)|
\pto 0.
\]
By \Cref{lem:topology-uniform},
\[
\sup_{w\in\Delta_J}
|L_{\mathrm{top},T_0}(w)-\mathcal L_{\mathrm{top}}(w)|
\pto 0.
\]
By Assumption A8,
\[
\widehat s_y \pto s_y,
\qquad
\widehat s_{\mathrm{top}}\pto s_{\mathrm{top}}.
\]
Since \(s_y,\widehat s_y\ge a>0\) and
\(s_{\mathrm{top}},\widehat s_{\mathrm{top}}\ge b>0\), standard ratio
arguments imply componentwise uniform convergence of the standardized losses.

To see this explicitly for the outcome loss,
\[
\begin{aligned}
\sup_{w\in\Delta_J}
\left|
\frac{L_{y,T_0}(w)}{\widehat s_y}
-
\frac{\mathcal L_y(w)}{s_y}
\right|
&\le
\frac{1}{\widehat s_y}
\sup_{w\in\Delta_J}
|L_{y,T_0}(w)-\mathcal L_y(w)|
\\
&\quad+
\left|
\frac{1}{\widehat s_y}
-
\frac{1}{s_y}
\right|
\sup_{w\in\Delta_J}\mathcal L_y(w).
\end{aligned}
\]
By Assumption A4,
\[
|Y_{1t}(0)-Y_{w,t}(0)|
\le
|Y_{1t}(0)|
+
\sum_{j=2}^{J+1}w_j|Y_{jt}(0)|
\le
2M,
\]
so
\[
\sup_{w\in\Delta_J}\mathcal L_y(w)\le4M^2.
\]
Thus the right-hand side above is \(o_p(1)\). The topology term is handled
identically using Assumptions A5 and A8.

Therefore,
\[
\begin{aligned}
&
\sup_{w\in\Delta_J}
\left|
\widetilde Q_{T_0,\lambda}(w)
-
\widetilde{\mathcal Q}_{\lambda}(w)
\right|
\\
&\le
\sup_{w\in\Delta_J}
\left|
\widetilde L_{y,T_0}(w)
-
\widetilde{\mathcal L}_y(w)
\right|
+
\lambda
\sup_{w\in\Delta_J}
\left|
\widetilde L_{\mathrm{top},T_0}(w)
-
\widetilde{\mathcal L}_{\mathrm{top}}(w)
\right|.
\end{aligned}
\]
For fixed \(\lambda<\infty\), the right-hand side is \(o_p(1)\). Hence
\[
\sup_{w\in\Delta_J}
\left|
\widetilde Q_{T_0,\lambda}(w)
-
\widetilde{\mathcal Q}_{\lambda}(w)
\right|
\pto 0.
\]
\end{proof}

% ------------------------------------------------------------
\subsection{Theorem 1: Uniform Consistency of TASCM}
% ------------------------------------------------------------

\begin{theorem}[Uniform consistency of TASCM]
\label{thm:consistency}
Suppose Assumptions A3--A9 hold. Suppose further that
\(\widehat w_\lambda\) is a \(\zeta_T\)-approximate minimizer with
\(\zeta_T\pto0\). Then, for fixed \(\lambda<\infty\),
\[
\widehat w_\lambda
\pto
w_\lambda^\star.
\]
\end{theorem}

\begin{proof}
The proof follows the standard extremum-estimator argument
\citep{newey1994large,vaart1998asymptotic}. Because \(w_\lambda^\star\) is the
unique minimizer of the continuous function \(\widetilde{\mathcal Q}_{\lambda}\)
on compact \(\Delta_J\), for every \(\epsilon>0\),
\[
\eta_\epsilon
=
\inf_{\|w-w_\lambda^\star\|\ge\epsilon}
\left[
\widetilde{\mathcal Q}_{\lambda}(w)
-
\widetilde{\mathcal Q}_{\lambda}(w_\lambda^\star)
\right]
>0.
\]
By \Cref{lem:objective-uniform} and \(\zeta_T\pto0\), with probability
approaching one,
\[
\sup_{w\in\Delta_J}
\left|
\widetilde Q_{T_0,\lambda}(w)
-
\widetilde{\mathcal Q}_{\lambda}(w)
\right|
<
\eta_\epsilon/4
\]
and
\[
\zeta_T<\eta_\epsilon/4.
\]

Suppose, toward contradiction, that
\[
\|\widehat w_\lambda-w_\lambda^\star\|\ge \epsilon.
\]
Then
\[
\widetilde{\mathcal Q}_{\lambda}(\widehat w_\lambda)
\ge
\widetilde{\mathcal Q}_{\lambda}(w_\lambda^\star)
+
\eta_\epsilon.
\]
Using uniform convergence,
\[
\widetilde Q_{T_0,\lambda}(\widehat w_\lambda)
\ge
\widetilde{\mathcal Q}_{\lambda}(w_\lambda^\star)
+
3\eta_\epsilon/4.
\]
Again using uniform convergence,
\[
\widetilde Q_{T_0,\lambda}(w_\lambda^\star)
\le
\widetilde{\mathcal Q}_{\lambda}(w_\lambda^\star)
+
\eta_\epsilon/4.
\]
Therefore,
\[
\widetilde Q_{T_0,\lambda}(\widehat w_\lambda)
>
\widetilde Q_{T_0,\lambda}(w_\lambda^\star)
+
\zeta_T,
\]
which contradicts the approximate-minimizer condition. Hence
\[
P(\|\widehat w_\lambda-w_\lambda^\star\|\ge \epsilon)\to0.
\]
Since \(\epsilon>0\) is arbitrary,
\[
\widehat w_\lambda\pto w_\lambda^\star.
\]
\end{proof}

% ------------------------------------------------------------
\subsection{Theorem 2: Oracle Inequality and Counterfactual Risk Bound}
% ------------------------------------------------------------

The next theorem corrects a common mistake in penalized-objective arguments.
From
\[
A+\lambda B
\le
\epsilon_y+\lambda\epsilon_{\mathrm{top}},
\]
one cannot conclude \(A\le \epsilon_y\) and
\(B\le \epsilon_{\mathrm{top}}\) separately. The correct conclusions are
\[
A\le \epsilon_y+\lambda\epsilon_{\mathrm{top}},
\]
and, if \(\lambda>0\),
\[
B\le \epsilon_y/\lambda+\epsilon_{\mathrm{top}}.
\]

\begin{theorem}[Oracle risk bound for TASCM]
\label{thm:oracle-risk}
Suppose Assumptions A10 and A11 hold. Suppose \Cref{lem:objective-uniform} holds
and \(\widehat w_\lambda\) is a \(\zeta_T\)-approximate minimizer with
\(\zeta_T\pto0\). Then, for any fixed \(\lambda>0\),
\[
\mathcal R(\widehat w_\lambda)
\le
C_y
\left[
\epsilon_{y,T}
+
\lambda\epsilon_{\mathrm{top},T}
\right]
+
C_{\mathrm{top}}
\left[
\frac{\epsilon_{y,T}}{\lambda}
+
\epsilon_{\mathrm{top},T}
\right]
+
r_T
+
o_p(1).
\]
Consequently, if
\[
\epsilon_{y,T}
+
\lambda\epsilon_{\mathrm{top},T}
\to0,
\]
\[
\frac{\epsilon_{y,T}}{\lambda}
+
\epsilon_{\mathrm{top},T}
\to0,
\]
and \(r_T\to0\), then
\[
\mathcal R(\widehat w_\lambda)\pto0.
\]
\end{theorem}

\begin{proof}
By approximate optimality,
\[
\widetilde Q_{T_0,\lambda}(\widehat w_\lambda)
\le
\widetilde Q_{T_0,\lambda}(w_T^\circ)
+
\zeta_T.
\]
Expanding gives
\[
\widetilde L_{y,T_0}(\widehat w_\lambda)
+
\lambda
\widetilde L_{\mathrm{top},T_0}(\widehat w_\lambda)
\le
\widetilde L_{y,T_0}(w_T^\circ)
+
\lambda
\widetilde L_{\mathrm{top},T_0}(w_T^\circ)
+
\zeta_T.
\]
By \Cref{lem:objective-uniform}, Assumption A11, and \(\zeta_T\pto0\),
\[
\widetilde L_{y,T_0}(\widehat w_\lambda)
+
\lambda
\widetilde L_{\mathrm{top},T_0}(\widehat w_\lambda)
\le
\epsilon_{y,T}
+
\lambda\epsilon_{\mathrm{top},T}
+
o_p(1).
\]
Since both losses are nonnegative,
\[
\widetilde L_{y,T_0}(\widehat w_\lambda)
\le
\epsilon_{y,T}
+
\lambda\epsilon_{\mathrm{top},T}
+
o_p(1),
\]
and, because \(\lambda>0\),
\[
\widetilde L_{\mathrm{top},T_0}(\widehat w_\lambda)
\le
\frac{\epsilon_{y,T}}{\lambda}
+
\epsilon_{\mathrm{top},T}
+
o_p(1).
\]
By componentwise uniform convergence,
\[
\widetilde{\mathcal L}_y(\widehat w_\lambda)
\le
\epsilon_{y,T}
+
\lambda\epsilon_{\mathrm{top},T}
+
o_p(1),
\]
and
\[
\widetilde{\mathcal L}_{\mathrm{top}}(\widehat w_\lambda)
\le
\frac{\epsilon_{y,T}}{\lambda}
+
\epsilon_{\mathrm{top},T}
+
o_p(1).
\]
Now apply Assumption A10:
\[
\mathcal R(\widehat w_\lambda)
\le
C_y\widetilde{\mathcal L}_y(\widehat w_\lambda)
+
C_{\mathrm{top}}\widetilde{\mathcal L}_{\mathrm{top}}(\widehat w_\lambda)
+
r_T.
\]
Substituting the two preceding bounds yields the stated inequality. If the two
bracketed terms and \(r_T\) converge to zero, then
\[
\mathcal R(\widehat w_\lambda)\pto0.
\]
\end{proof}

\begin{remark}[Rate-balancing role of \(\lambda\)]
The oracle bound displays a tradeoff in \(\lambda\). Ignoring constants, the
\(\lambda\)-dependent terms are
\[
\lambda\epsilon_{\mathrm{top},T}
+
\epsilon_{y,T}/\lambda.
\]
If the oracle rates were known, a rate-balancing choice would satisfy
\[
\lambda_T
\asymp
\sqrt{
\frac{\epsilon_{y,T}}{\epsilon_{\mathrm{top},T}}
}.
\]
The validation rule below provides a feasible data-driven alternative within a
pre-specified grid.
\end{remark}

\begin{remark}[Varying \(\lambda_T\)]
If \(\lambda=\lambda_T\) is allowed to vary with \(T_0\), the same argument
holds provided the uniform approximation errors are small relative to
\(\lambda_T\). In particular, if the uniform convergence rate is
\(\rho_T=o_p(1)\), then the topological component contains a term of order
\(\rho_T/\lambda_T\). Thus one needs \(\rho_T/\lambda_T=o_p(1)\) in addition to
the displayed conditions.
\end{remark}

% ------------------------------------------------------------
\subsection{Theorem 3A: Validation Optimality for \texorpdfstring{\(\lambda\)}{lambda}}
% ------------------------------------------------------------

Let
\[
\Lambda=\{\lambda_1,\ldots,\lambda_K\}
\]
be a fixed finite grid containing \(0\). The grid is fixed before observing
post-treatment outcomes. Split the pre-treatment period into a training block
\(\mathcal T_{\train}\) and a validation block \(\mathcal T_{\val}\). We assume
that both blocks contain more than \((m-1)\ell\) observations, so that delay
embeddings and topology losses are well-defined on both blocks.

For each \(\lambda\in\Lambda\), let
\(\widehat w_\lambda^{\train}\) be the TASCM weight estimated using only the
training block. Define validation outcome fit
\[
R_T(\lambda)
=
\RMSPE_{\val}(\lambda),
\]
and validation topology loss
\[
G_T(\lambda)
=
L_{\mathrm{top},\val}(\lambda).
\]
Let \(R(\lambda)\) and \(G(\lambda)\) denote their population validation
counterparts. For \(c>1\) and a fixed floor \(a_R>0\), define the empirical
admissible set
\[
\widehat\Lambda_{\adm}
=
\left\{
\lambda\in\Lambda:
R_T(\lambda)
\le
c\left(R_T(0)\vee a_R\right)
\right\}.
\]
Define the population admissible set
\[
\Lambda_{\adm}^0
=
\left\{
\lambda\in\Lambda:
R(\lambda)
\le
c\left(R(0)\vee a_R\right)
\right\}.
\]
Because \(0\in\Lambda\), the empirical admissible set is nonempty. The
lexicographic validation rule is
\[
\widehat\lambda
=
\min
\argmin_{\lambda\in\widehat\Lambda_{\adm}}
G_T(\lambda).
\]
After selecting \(\widehat\lambda\), the final TASCM weights are re-estimated on
the full pre-treatment period using \(\widehat\lambda\). The validation split is
used only for selecting \(\lambda\).

\begin{theorem}[Lexicographic validation optimality]
\label{thm:validation}
Assume \(\Lambda\) is finite and
\[
\sup_{\lambda\in\Lambda}
|R_T(\lambda)-R(\lambda)|
\pto0,
\]
\[
\sup_{\lambda\in\Lambda}
|G_T(\lambda)-G(\lambda)|
\pto0.
\]
Assume further that no candidate lies on the population admissibility boundary:
\[
R(\lambda)
\neq
c\left(R(0)\vee a_R\right)
\qquad
\text{for all }\lambda\in\Lambda,
\]
and that
\[
\lambda_{\val}^{\star}
=
\min
\argmin_{\lambda\in\Lambda_{\adm}^0}
G(\lambda)
\]
is unique. Then
\[
\widehat\lambda \pto \lambda_{\val}^{\star},
\]
and
\[
G(\widehat\lambda)
\le
\inf_{\lambda\in\Lambda_{\adm}^0}G(\lambda)
+
o_p(1).
\]
\end{theorem}

\begin{proof}
Since \(\Lambda\) is finite, pointwise convergence implies uniform convergence.
By the no-boundary condition, the empirical admissible set equals the population
admissible set with probability approaching one:
\[
P(\widehat\Lambda_{\adm}=\Lambda_{\adm}^0)\to1.
\]
On this event,
\[
\widehat\lambda
=
\min
\argmin_{\lambda\in\Lambda_{\adm}^0}
G_T(\lambda).
\]
Uniform convergence of \(G_T\) to \(G\) over the finite set
\(\Lambda_{\adm}^0\) and uniqueness of the minimizer imply finite-grid argmin
consistency:
\[
\widehat\lambda \pto \lambda_{\val}^{\star}.
\]
Furthermore,
\[
G(\widehat\lambda)
=
G_T(\widehat\lambda)+o_p(1)
\le
G_T(\lambda_{\val}^{\star})+o_p(1)
=
G(\lambda_{\val}^{\star})+o_p(1).
\]
Since
\[
G(\lambda_{\val}^{\star})
=
\inf_{\lambda\in\Lambda_{\adm}^0}G(\lambda),
\]
we obtain
\[
G(\widehat\lambda)
\le
\inf_{\lambda\in\Lambda_{\adm}^0}G(\lambda)+o_p(1).
\]
\end{proof}

% ------------------------------------------------------------
\subsection{Corollary 3B: A Sufficient Condition for Post-Treatment Risk Consistency}
% ------------------------------------------------------------

The preceding theorem proves validation optimality, not post-treatment risk
optimality. To convert validation optimality into post-treatment risk
consistency, one needs an additional representativeness condition.

\begin{corollary}[Sufficient risk consistency condition]
\label{cor:risk-consistency}
Suppose \Cref{thm:validation} holds. Suppose further that, uniformly over
\(\Lambda_{\adm}^0\),
\[
\left|
G(\lambda)-\mathcal R(w_\lambda^\star)
\right|
\to0.
\]
Then
\[
\mathcal R(w_{\widehat\lambda}^\star)
\le
\inf_{\lambda\in\Lambda_{\adm}^0}
\mathcal R(w_\lambda^\star)
+
o_p(1).
\]
\end{corollary}

\begin{proof}
By \Cref{thm:validation},
\[
G(\widehat\lambda)
\le
\inf_{\lambda\in\Lambda_{\adm}^0}G(\lambda)+o_p(1).
\]
By the validation-risk representativeness assumption,
\[
\mathcal R(w_{\widehat\lambda}^\star)
=
G(\widehat\lambda)+o_p(1),
\]
and uniformly over \(\lambda\in\Lambda_{\adm}^0\),
\[
G(\lambda)
=
\mathcal R(w_\lambda^\star)+o(1).
\]
Therefore,
\[
\mathcal R(w_{\widehat\lambda}^\star)
\le
\inf_{\lambda\in\Lambda_{\adm}^0}
\mathcal R(w_\lambda^\star)
+
o_p(1).
\]
\end{proof}

\paragraph{Important qualification.}
\Cref{thm:validation} is the main validation theorem. \Cref{cor:risk-consistency}
requires the stronger assumption that validation topology loss is representative
of post-treatment counterfactual risk. This condition is sufficient but not
necessary. It should not be interpreted as a general equivalence between
validation topology loss and post-treatment counterfactual risk.

% ------------------------------------------------------------
\subsection{A Pairwise Sufficient Condition for Topology Dominance}
% ------------------------------------------------------------

This theorem characterizes when topology can overturn the ranking produced by
outcome-only SCM. It is a pairwise comparison theorem, not a general dominance
theorem.

Let \(w_y\) be a level-best synthetic control:
\[
w_y
\in
\argmin_{w\in\Delta_J}\widetilde{\mathcal L}_y(w).
\]
If the level-best minimizer is not unique, fix any selected minimizer \(w_y\).
Let \(w_g\in\Delta_J\) be a topology-matching candidate. Suppose
\[
\widetilde{\mathcal L}_y(w_y)
<
\widetilde{\mathcal L}_y(w_g),
\]
so that outcome-only SCM prefers \(w_y\). Suppose also
\[
\widetilde{\mathcal L}_{\mathrm{top}}(w_g)
<
\widetilde{\mathcal L}_{\mathrm{top}}(w_y),
\]
so that \(w_g\) is topologically closer to the treated unit. Finally suppose
\[
\mathcal R(w_g)<\mathcal R(w_y),
\]
so that \(w_g\) is genuinely better for post-treatment counterfactual risk.

Define
\[
\lambda_{\mathrm{dom}}
=
\frac{
\widetilde{\mathcal L}_y(w_g)
-
\widetilde{\mathcal L}_y(w_y)
}{
\widetilde{\mathcal L}_{\mathrm{top}}(w_y)
-
\widetilde{\mathcal L}_{\mathrm{top}}(w_g)
}.
\]
The numerator and denominator are positive, so \(\lambda_{\mathrm{dom}}>0\).

\begin{theorem}[Pairwise topology-dominance region]
\label{thm:pairwise-dominance}
Under the three inequalities above, for every
\[
\lambda>\lambda_{\mathrm{dom}},
\]
we have
\[
\widetilde{\mathcal L}_y(w_g)
+
\lambda\widetilde{\mathcal L}_{\mathrm{top}}(w_g)
<
\widetilde{\mathcal L}_y(w_y)
+
\lambda\widetilde{\mathcal L}_{\mathrm{top}}(w_y).
\]
Thus, for all \(\lambda>\lambda_{\mathrm{dom}}\), the TASCM objective ranks the
topology-matching candidate \(w_g\) above the level-best SCM candidate \(w_y\).
\end{theorem}

\begin{proof}
We need to show that
\[
\widetilde{\mathcal L}_y(w_g)
+
\lambda\widetilde{\mathcal L}_{\mathrm{top}}(w_g)
<
\widetilde{\mathcal L}_y(w_y)
+
\lambda\widetilde{\mathcal L}_{\mathrm{top}}(w_y).
\]
Rearranging gives
\[
\widetilde{\mathcal L}_y(w_g)
-
\widetilde{\mathcal L}_y(w_y)
<
\lambda
\left[
\widetilde{\mathcal L}_{\mathrm{top}}(w_y)
-
\widetilde{\mathcal L}_{\mathrm{top}}(w_g)
\right].
\]
Because
\[
\widetilde{\mathcal L}_{\mathrm{top}}(w_y)
-
\widetilde{\mathcal L}_{\mathrm{top}}(w_g)
>0,
\]
this inequality is equivalent to
\[
\lambda>\lambda_{\mathrm{dom}}.
\]
Therefore, for every \(\lambda>\lambda_{\mathrm{dom}}\), the TASCM objective
ranks \(w_g\) above \(w_y\).
\end{proof}

% ------------------------------------------------------------
\subsection{Corollary 4B: Global Dominance in an Admissible Region}
% ------------------------------------------------------------

The previous theorem is pairwise. It does not imply that TASCM globally selects
\(w_g\). To obtain a global dominance result, impose an admissible-fit region
\[
\mathcal A_c
=
\left\{
w\in\Delta_J:
\widetilde{\mathcal L}_y(w)
\le
c\widetilde{\mathcal L}_y(w_y)
\right\},
\qquad c>1.
\]

\begin{corollary}[Global topology dominance under admissible minimization]
\label{cor:global-dominance}
Suppose the conditions of \Cref{thm:pairwise-dominance} hold. Suppose
additionally that, for some \(\lambda>\lambda_{\mathrm{dom}}\),
\[
w_g
\in
\argmin_{w\in\mathcal A_c}
\left\{
\widetilde{\mathcal L}_y(w)
+
\lambda\widetilde{\mathcal L}_{\mathrm{top}}(w)
\right\}.
\]
Then TASCM restricted to the admissible region selects \(w_g\). Moreover,
\[
\mathcal R(w_g)<\mathcal R(w_y).
\]
Hence, in this admissible region, topology-augmented SCM achieves lower
post-treatment counterfactual risk than the selected level-best SCM control.
\end{corollary}

\begin{proof}
By assumption, \(w_g\) is a global minimizer of the TASCM objective over
\(\mathcal A_c\). Therefore the admissible TASCM estimator selects \(w_g\).
The strict risk improvement,
\[
\mathcal R(w_g)<\mathcal R(w_y),
\]
is one of the maintained pairwise dominance conditions in
\Cref{thm:pairwise-dominance}. Hence the selected topology-augmented control
has lower post-treatment counterfactual risk than the selected level-only SCM
candidate \(w_y\).
\end{proof}

% ------------------------------------------------------------
\subsection{Constructive Cyclic DGP: Sufficient Conditions for Pairwise Topology Dominance}
% ------------------------------------------------------------

This subsection gives a sufficient-condition example showing when the pairwise
dominance inequalities can arise. It is not a universal dominance result.

Consider the untreated outcome process
\[
Y_{it}(0)
=
\alpha_i+\beta_i t/T_0
+
A_i\sin(\omega_i t+\phi_i)
+
\varepsilon_{it},
\]
where \(\varepsilon_{it}\) is mean zero with uniformly bounded second moment.
The treated unit satisfies
\[
A_1=A^*,
\qquad
\omega_1=\omega^*,
\qquad
\phi_1=\phi^*.
\]

Suppose there are two candidate synthetic controls. The first, \(w_y\), is
level-fitting but topology-wrong. It matches the trend component well but does
not match the cyclic component. The second, \(w_g\), is topology-correct. It has
slightly worse pre-treatment level fit but matches the cyclic component and
therefore the standardized delay-embedding loop structure.

Assume there exist constants
\[
\Delta_y>0,\qquad
\Delta_{\mathrm{top}}>0,\qquad
\Delta_R>0
\]
such that
\[
\widetilde{\mathcal L}_y(w_g)
-
\widetilde{\mathcal L}_y(w_y)
=
\Delta_y,
\]
\[
\widetilde{\mathcal L}_{\mathrm{top}}(w_y)
-
\widetilde{\mathcal L}_{\mathrm{top}}(w_g)
=
\Delta_{\mathrm{top}},
\]
and
\[
\mathcal R(w_y)-\mathcal R(w_g)
=
\Delta_R.
\]
Then
\[
\lambda_{\mathrm{dom}}
=
\frac{\Delta_y}{\Delta_{\mathrm{top}}}.
\]
For every \(\lambda>\lambda_{\mathrm{dom}}\), the pairwise topology-dominance
result implies that TASCM ranks \(w_g\) above \(w_y\). If \(w_g\) is the
minimizer in the admissible region, the global dominance corollary implies that
TASCM has lower post-treatment risk.

\paragraph{Comment.}
A fully closed-form proof of the gap
\(\Delta_{\mathrm{top}}>0\) would require lower bounds on the persistence
diagram distance between standardized delay embeddings of matched and
mismatched sinusoidal trajectories. The present result is stated as a
sufficient-condition theorem: when a cyclic DGP generates a candidate with
slightly worse level fit but strictly better topological fit and post-treatment
risk, TASCM improves upon level-only SCM for
\(\lambda>\Delta_y/\Delta_{\mathrm{top}}\).

% ------------------------------------------------------------
\subsection{Vietoris--Rips Implementation Corollary}
% ------------------------------------------------------------

The main population theory is stated for a generic stable filtration
\(\mathcal F\). The empirical implementation may use Vietoris--Rips persistence
on finite standardized delay-embedded trajectories. This requires stronger
assumptions.

Let
\[
E_{m,\ell}^{S}(Y_w^{\pre})
=
\left\{
X_{w,t}^{S,m,\ell}:
t=(m-1)\ell+1,\ldots,T_0
\right\}
\]
be the finite standardized delay-embedded point cloud. Let
\(\mathcal S_w^{S}\) denote the population support of the standardized
delay-window law \(P_w^{S}\). Suppose
\[
\sup_{w\in\Delta_J}
d_H
\left(
E_{m,\ell}^{S}(Y_w^{\pre}),
\mathcal S_w^{S}
\right)
\pto 0.
\]
Assume Vietoris--Rips persistence is stable in Hausdorff distance and \(\Phi\)
is Lipschitz with respect to the persistence-diagram metric. Then
\[
\sup_{w\in\Delta_J}
\left|
L_{\mathrm{top},T_0}^{S,\mathrm{VR}}(w)
-
\mathcal L_{\mathrm{top}}^{S,\mathrm{VR}}(w)
\right|
\pto 0.
\]
Therefore, the M-estimation consistency result in \Cref{thm:consistency} also
applies to the VR-based finite-sample TASCM implementation under these stronger
Hausdorff/support-convergence assumptions.

\paragraph{Interpretation.}
DTM or kernel-distance filtrations are cleaner for population theory because
they are directly defined from probability laws and have distributional
stability. Vietoris--Rips persistence is natural for finite-sample
implementation and visualization, but it relies on stronger support-level
conditions and is more sensitive to outliers. This Hausdorff support-convergence
condition is stronger than the weak or distributional convergence required in
Assumption A6 and is not implied by A6 in general.

% ------------------------------------------------------------
\subsection{Implication for Treatment-Effect Estimation and Inference}
% ------------------------------------------------------------

For \(t>T_0\),
\[
\widehat\tau_{1t}^{\TASCM}-\tau_{1t}
=
\left[
Y_{1t}-\widehat Y_{1t}^{\TASCM}(0)
\right]
-
\left[
Y_{1t}(1)-Y_{1t}(0)
\right].
\]
Since \(Y_{1t}=Y_{1t}(1)\) for the treated unit after treatment,
\[
\widehat\tau_{1t}^{\TASCM}-\tau_{1t}
=
Y_{1t}(0)
-
\sum_{j=2}^{J+1}\widehat w_jY_{jt}(0).
\]
Thus treatment-effect estimation error is exactly counterfactual approximation
error.

Analytic standard errors are not immediate because TASCM contains a nonlinear
persistent-homology term and the optimizer may be non-smooth in the data. In
applications, uncertainty can be assessed using placebo or permutation
inference, re-estimating the TASCM procedure for each placebo unit. One may
either re-select \(\lambda\) within each placebo iteration or condition on the
selected \(\widehat\lambda\). The former corresponds to procedure-level
inference, while the latter is conditional on the selected tuning parameter.

% ------------------------------------------------------------
\subsection{Summary of Corrected Results}
% ------------------------------------------------------------

The corrected theoretical package is as follows.

First, TASCM is a standardized approximate M-estimator. For a fixed
\(\lambda\ge0\), the empirical objective is
\[
\widetilde Q_{T_0,\lambda}(w)
=
\widetilde L_{y,T_0}(w)
+
\lambda
\widetilde L_{\mathrm{top},T_0}(w).
\]
Here \(L_{\mathrm{top},T_0}(w)\) is computed from the standardized synthetic
trajectory \(S_{T_0}(Y_w^{\pre})\), after forming
\[
Y_w^{\pre}
=
\sum_{j=2}^{J+1}w_jY_j^{\pre}.
\]
The TASCM estimator \(\widehat w_\lambda\) is defined as a
\(\zeta_T\)-approximate minimizer satisfying
\[
\widetilde Q_{T_0,\lambda}(\widehat w_\lambda)
\le
\inf_{w\in\Delta_J}
\widetilde Q_{T_0,\lambda}(w)
+
\zeta_T,
\qquad
\zeta_T\pto0.
\]
The exact global-minimizer case corresponds to \(\zeta_T=0\).

Its population target is
\[
w_\lambda^\star
\in
\argmin_{w\in\Delta_J}
\left\{
\widetilde{\mathcal L}_y(w)
+
\lambda
\widetilde{\mathcal L}_{\mathrm{top}}(w)
\right\}.
\]

Second, under uniform convergence, compactness, uniqueness of the population
minimizer, and \(\zeta_T\pto0\),
\[
\widehat w_\lambda \pto w_\lambda^\star.
\]

Third, under the approximate topology-augmented representation condition and
the pre-post stability condition, the penalized oracle risk bound is
\[
\mathcal R(\widehat w_\lambda)
\le
C_y
\left[
\epsilon_{y,T}
+
\lambda\epsilon_{\mathrm{top},T}
\right]
+
C_{\mathrm{top}}
\left[
\frac{\epsilon_{y,T}}{\lambda}
+
\epsilon_{\mathrm{top},T}
\right]
+
r_T
+
o_p(1).
\]
This is the valid risk-control statement. It does not imply separate
component-wise oracle bounds unless one uses a constrained estimator.

Fourth, blocked validation yields validation optimality:
\[
G(\widehat\lambda)
\le
\inf_{\lambda\in\Lambda_{\mathrm{adm}}^0}
G(\lambda)
+
o_p(1).
\]
Post-treatment risk consistency of \(\widehat\lambda\) requires the additional
representativeness condition stated in \Cref{cor:risk-consistency}.

Fifth, topology dominance is pairwise and conditional. If
\[
\lambda
>
\lambda_{\mathrm{dom}}
=
\frac{
\widetilde{\mathcal L}_y(w_g)
-
\widetilde{\mathcal L}_y(w_y)
}{
\widetilde{\mathcal L}_{\mathrm{top}}(w_y)
-
\widetilde{\mathcal L}_{\mathrm{top}}(w_g)
},
\]
then TASCM ranks \(w_g\) above \(w_y\). Global dominance follows only under the
additional admissible-region minimization condition in
\Cref{cor:global-dominance}.

\bibliographystyle{apalike}
\bibliography{references}

\end{document}